%% ============================================================
%% 分片 34：原文第 11 节「$A$-Branes and ${\cal D}$-Modules」（原文 glshort.tex 起于第 9938 行）
%% 由 tools/harvmac2tex.py 从 arXiv 源码（hep-th/0604151v3）自动生成。
%% 只译英文散文；公式、编号、\label、\cite 一律不改。
%% ============================================================

\subsection{典范余各向同性（c.c.）膜的推广与扭化微分算子}\label{sec:genotics}

我们已经看到，几何 Langlands 对偶可以理解为交换 $\Psi=0$ 与
$\Psi=\infty$ 的 $S$-变换的结果。更一般地，变换 $S$ 把规范群为
$G$、典范参数为 $\Psi$ 的规范场论映为规范群为 $\LG$、典范参数为
$-1/n_{\frak g}\Psi$ 的规范场论（其中，正如我们在
第~\ref{sec:reviewsd}~节所解释的，若 $G$ 是单连通型，则
$n_{\frak g}=1$，否则 $n_{\frak
g}=2$ 或 $3$）。我们在此的目标是描述：当我们偏离
$\Psi=0,\infty$ 时，会出现何种形式的
几何 Langlands 纲领的推广。（关于经由共形场论的这一推广的讨论，
见 \cite{ref9} 第 6.3 与 8.6 节；阿贝尔情形见 \cite{ref146}。）

我们将要探讨的几何 Langlands 纲领的这种推广，或许在 $\Psi$ 取
有理数值时最为有趣，因为正如我们在第~\ref{sec:tophooft}~节所见，
正是在这一情形才存在有意思的线算子。此外，完整的
$S$-对偶群——正如我们在第~\ref{sec:canonpar}~节所证明的，它由
$T:\Psi\to\Psi+1$ 连同 $S:\Psi\to -1/n_{\frak
g}\Psi$ 生成——把 $\Psi$ 的有理值映为有理值。事实上，$\Psi$
的一切有理值都被这个群联系在一起；这一点也保证了在 $\Psi$ 取
有理值时存在大量膜，正如在 $\Psi=0$ 与 $\Psi=\infty$ 时那样。

不过，这一对偶对一切复值的 $\Psi$ 都成立。但为简单起见，这里
我们取 $\Psi$ 为实数。只要令 $t=1$ 并让 $\theta$ 角任意，就能得到
$\Psi$ 的任意实数值： % 式 (11.26)
\setcounter{equation}{25}
\begin{equation}\label{eq:melo}
\Psi={\theta\over
{2\pi}}=\Re\tau.
\end{equation}
 $S$-对偶一般会变换 ${\rm
Re}\,\tau=\theta$、${\rm Im}\,\tau=4\pi/e^2$ 以及 $t$。但 $\Psi$
才是唯一重要的参数，所以只要 $\Psi$ 是实数，
我们就可以把 $S$-对偶看作只作用在 $\theta$ 上。在这一设置下，
「电」理论与「磁」理论看起来对称得多：两者都有 $t=1$，但它们的
$\theta$ 角互为反比。在黎曼面 $C$ 上紧化之后，便得到复结构 $K$
中以 $\MH(G,C)$ 为靶的 $A$-模型，以及 $B$-场 % 式 (11.27)
\setcounter{equation}{26}
\begin{equation}\label{eq:pelbo}
B=-\omega_I \Re\tau,
\end{equation}
 正如式 \eqref{eq:palooka} 所说明的。$S$-变换的实现方式是
把 $\Re\tau$ 换成 $-1/n_{\frak g}\Re\tau$、把 $G$ 换成
$^L\negthinspace G$。

由 \eqref{eq:pelbo}\ 给出的 $B$-场沿 Hitchin 纤维化的纤维并不平坦；
事实上，它是复结构 $I$ 中 Kähler 形式的一个倍数，因而限制到
一条纤维上是非退化的。一个重要的推论是：对一般的 $\Psi$，
Hitchin 纤维化的纤维不再是有效的 $A$-膜。为解释这一点，
回忆一下：靶为 $X$ 的 sigma 模型在存在一个包裹于子流形 $Y$ 上的
$D$-膜时，在如下规范变换下保持不变 % 式 (11.28)
\setcounter{equation}{27}
\begin{equation}\label{eq:zebol}
B\to
B-d\lambda,\quad A\to A+\lambda\vert_Y,\quad \lambda\in
\Omega^1(X).
\end{equation}
 规范不变的组合是 $F+B\vert_Y$。现在
假设 $X$ 是辛流形、$Y$ 是 Lagrangian 子流形。若 $B=0$，则包裹在
$Y$ 上的单个 $D$-膜只有当 $F=0$ 时才能是 $A$-膜。
这一条件的规范不变推广（对 $B\neq 0$ 也成立）是 % 式 (11.29)
\setcounter{equation}{28}
\begin{equation}\label{eq:subol}
F+B\vert_Y=0.
\end{equation}
 若
$B|_Y/2\pi$ 的上同调类不是整的，那么这个方程对任何 $U(1)$ 规范场
都无法满足，因此单个 $A$-膜不能包裹在 $Y$ 上。在我们的
情形，这意味着当 $\theta\neq 2\pi n$ 时，Hitchin 纤维化的纤维不是
有效的 $A$-膜，无论怎样选取 Chan-Paton 丛。

因此，当 $\Psi$ 非整数时，研究对象将不是支撑在 Hitchin 纤维化的
一条纤维上的秩一膜。
 作为一致性检验，注意它们朴素的 $T$-对偶，即
$\MH(\LG,C)$ 上的点，在 $t=1$ 且 $\theta\neq 0$ 时也不是有效的
拓扑膜。事实上，点永远不是 $A$-膜，因为 $A$-膜的支撑总是至少
中间维数的。当 $\Psi$ 为有理数时，存在支撑在 Hitchin 纤维化的
一条纤维上的有限秩 $A$-膜；我们可以从 $\Psi=0$ 出发沿对偶走一遍，
找出它们是什么。
无论如何，$S$-对偶仍然把 $\MH(G,C)$ 的 $A$-膜映为
$\MH(\LG,C)$ 的 $A$-膜，尽管 $A$-膜的例子可能很稀少。

但当 $\Psi\not=0$ 时，关于 $A$-膜我们能说些什么，才能推广我们在
$\Psi=0$ 时得到的结果呢？在本小节的剩余部分，我们将回答这个问题。
 我们将在经典极限 $e^2\to
0,\ \Im\,\tau\to \infty$ 附近做全部计算，但由于理论只依赖于
$\Psi=\Re\tau$，所得结果对一切 $e^2$ 值都成立。

我们将看到，对一切 $\Psi$ 都存在典范余各向同性（c.c.）膜的类比物。
我们将把这个膜称为 ${\cal B}_{c.c.}^\Psi$。它的开弦代数
将是 ${\cal
D}_{K_{\EUM}^{1/2}\otimes {{{\frak L}}}^\Psi}$，也就是
$K_\M^{1/2}\otimes{\frak
L}^\Psi$ 上全纯微分算子的层，其中 ${\frak L}$ 是
$\cM(G,C)$ 上的行列式线丛（见对式 \eqref{eq:frox} 的讨论）。给定
$\MH(G,C)$ 上任何其他 $A$-膜 ${\cal B}'$，$({\cal
B}_{c.c.}^\Psi,{\cal B}')$ 弦的层是 ${\cal D}_{K_{
\EUM}^{1/2}\otimes {{{\frak L}}}^\Psi}$ 的一个模。因此，一般 $\Psi$
下的一个 $A$-膜给出此代数的一个模层。正如在
$\Psi=0$ 时一样，这一构造实际上来自四维
规范场论；用代数几何的话说，这意味着这些扭化
${\cal D}$-模可以延拓到所有
$G$-丛构成的叠上。$S$-对偶于是把 $G$ 的扭化 ${\cal D}$-模
映为 $^L\negthinspace G$ 的相应对象，同时作替换 $\Psi\to
-1/n_{\frak g}\Psi$。

充满空间的 $A$-膜上的规范场必须解方程 % 式 (11.30)
\setcounter{equation}{29}
\begin{equation}\label{eq:juko}
(\omega^{-1}(F+B))^2=-1,
\end{equation}
 其中
$\omega=(\Im\tau)\, \omega_K$。一个显然的解是 % 式 (11.31)
\setcounter{equation}{30}
\begin{equation}\label{eq:uko}
F=
\Im\tau\, (\cos{\eurm q})\, \omega_J,\quad \sin{\eurm q}=-{{\Re\,
\tau}\over {\Im\, \tau}}.
\end{equation}
  我们取这个解来定义膜 ${\cal B}_{c.c.}^\Psi$。注意，由于
${\Im\,\tau}$ 可以任意大，${\eurm q}$ 极其小，最终有意思的效应
都将线性地依赖于 $\Psi=\Re\,\tau$，而它是 ${\eurm q}$ 阶的量。

如果忽略代数结构，那么开弦层局部同构于复结构\footnote{与式 \eqref{eq:wyre} 比较，
我们看到 $I(\Psi)$ 与 $w=i\tan({\eurm
q}/2)$ 时的 $I_w$ 相同。} % 式 (11.32)
\setcounter{equation}{31}
\begin{equation}\label{eq:muko}
I(\Psi)=\omega^{-1}(F+B)=I \cos{\eurm q} -
J\sin{\eurm q}.
\end{equation}
 下 $\MH(G,C)$ 上全纯函数的层。相应的全纯坐标
是 % 式 (11.33)
\setcounter{equation}{32}
\begin{equation}\label{eq:suko}
A'_\bz=A_\bz - i \tan{{\eurm q}\over
2}\, \phi_\bz,\quad \phi'_z=\phi_z+i\tan{{\eurm q}\over 2}\, A_z.
\end{equation}

照例，$A$ 是 $G$-丛 $E\to C$ 上的联络，$\phi$ 是取值于
${\rm ad}(E)$ 的一形式。我们在第~\ref{sec:hitchfib}~节所称的
Hitchin 第二纤维化，就是遗忘映射
% 式 (11.34)
\setcounter{equation}{33}
\begin{equation}\label{eq:xnonko}
\pi_0: (A_\bz,\phi_z)\to A_\bz
\end{equation}
 从 $\MH$ 到 $\M$。显然，除 $\Psi=0$ 外，$\pi_0$ 关于 $I(\Psi)$
不是全纯的。为纠正这一状况，我们在 $E$ 上定义一个新的
$\bpartial$-联络 % 式 (11.35)
\setcounter{equation}{34}
\begin{equation}\label{eq:hoxo}
D'_\bz=D_\bz-i
\tan{{\eurm q}\over 2}\,\phi_\bz=\partial_\bz+A_\bz - i
\tan{{\eurm q}\over 2}\,\phi_\bz,
\end{equation}
 并把对 $(A,\phi)$ 映入由 $D'_\bz$ 定义的 $C$ 上全纯丛。让我们
把这个映射记为 $\pi_\Psi$；它显然关于 $I(\Psi)$ 全纯，且在
$\Psi=0$ 时约化为 $\pi_0$。照例，我们把这一纤维化限制到如下的
Higgs 丛子集上：$E$ 配以修改后的 $\bar\partial$ 算子后是稳定的。

映射 $\pi_\Psi$ 的纤维是什么样子？为理解这一点，注意纤维坐标的
显然候选者 $\phi_z+i\tan{{\eurm q}\over 2}\, A_z$ 在规范变换下非齐次
地变换。齐次变换的对象是 % 式 (11.36)
\setcounter{equation}{35}
\begin{equation}\label{eq:gyko}
D'_z=\phi_z+i\tan{{\eurm
q}\over 2}\, D_z,
\end{equation}
 其中 $D_z$ 是 $E$ 上的全纯协变导数。差一个因子
$i\tan{{\eurm q}\over 2}$，$D'_z$ 是丛 $E$ 上的一个
$\partial$-联络。我们选择不除以 $i\tan{{\eurm q}\over 2}$，
因为我们希望这些公式即使在 $\Psi=0$ 时也有意义。于是 $D'_z$
就是所谓 $E$ 上的全纯 $\lambda$-联络，其中
$\lambda=i\tan{{\eurm q}\over 2}$。对任何复流形 $X$，全纯丛
$E\to X$ 上的全纯 $\lambda$-联络是一个线性映射
$\nabla: \Gamma(E)\ra \Gamma(E\ot T^*X)$，它与 $E$ 的
$\bar\partial$ 算子交换，并且对任何函数 $f$ 与 $E$ 的任何截面
$s$ 满足
% 式 (11.37)
\setcounter{equation}{36}
\begin{equation}\label{eq:hyfo}
\nabla (f\cdot s)=f\nabla s + \lambda
\partial f\wedge s,
\end{equation}
 当 $\lambda=1$ 时，$\lambda$-联络就是 $E$ 上普通的
$\partial$ 算子；当 $\lambda=0$ 时，它是 $\End(E)\ot T^*X$ 的
一个截面。

借助 Hitchin 方程，$D'_z$ 与 $D'_\bz$ 交换（这与
第~\ref{sec:complexstr}~节末尾关于复结构 $I_w$ 的讨论密切相关），
因此定义了全纯丛 $E'=\pi_\Psi(E,\phi)$ 上的全纯
$\lambda$-联络。按通常的逻辑，这意味着 $E'$ 上的纤维是全纯
$\lambda$-联络的模空间。

固定向量丛 $E'\to C$ 上全纯 $\lambda$-联络的空间，自然地是一个
仿射空间，它以 $\End(E')\ot K_C$ 的全纯截面空间为模型空间，其中
$K_C$ 是 $C$ 的典范线丛。（这句话的意思是：我们可以给一个全纯
$\lambda$-联络加上 $\End(E')\otimes K_C$ 的一个全纯截面，而任何
两个全纯 $\lambda$-联络都以这种方式相联系。）空间
$H^0(C,\End(E')\ot K_C)$ 可以等同于 $\cM(G,C)$ 在点 $E'$ 处的
余切空间。于是在复结构 $I(\Psi)$ 中，模空间 $\MH(G,C)$ 看起来
像是 $T^*\cM(G,C)$ 的仿射形变；换句话说，它是仿射空间的丛，其
底层的向量丛是 $T^*\CM(G,C)$。总之，如果忽略代数结构，
$({\cal B}_{c.c.}^\Psi,{\cal B}_{c.c.}^\Psi)$ 弦的层是
$\cM(G,C)$ 的一个「扭化」余切丛上全纯函数的层。

剩下的问题是确定 $({\cal B}_{c.c}^\Psi,{\cal B}_{c.c.}^\Psi)$ 弦的
层的代数结构。我们把自己限制在沿仿射丛的纤维至多多项式增长的
那些函数上。\footnote{这一限制相当于在赋以复结构 $I(\Psi)$ 的
$\MH$ 上选取一个特定的代数结构。在这一代数结构中，它被称为带有
$\lambda$-联络的 $G$-丛的模空间。还有一个与同一复结构相容的
第二种代数结构，即在其中 $\MH$ 等价于稳定同态
$\pi_1(C)\to G_{\Bbb{C}}$ 的模空间。在这第二种代数结构里，
holonomy 的迹被当作代数函数。我们刚刚遇到的情形，是本文中唯一
一处必须在与 $\MH$ 上固定复结构相容的不同代数结构之间作选择的
地方。这一选择在几何 Langlands 纲领的其他进路中扮演着更突出的
角色。} 我们可以遵循与 $\Psi=0$ 时相同的逻辑。
 局部地说，我们量子化的是扭化余切丛而不是普通余切丛，这一点
无关紧要。整体地说，它可能会（而且正如我们将看到的，确实会）
修改 \eqref{eq:hecox} 中的转移函数。于是在任何 $\Psi$ 下，c.c.
弦的层都是作用在 $\cM(G,C)$ 上某个全纯线丛的某个幂上的全纯
微分算子层。

若 $\Psi$ 是整数，我们可以说得更精确。把 $\Psi$ 平移
$\Psi+n$ 会把 $B$-场——更确切地说 $B/2\pi$——平移线丛
${{{\frak L}}}^{-n}$ 的第一 Chern 类。（我们可以从
\eqref{eq:palooka} 看出这一点；平移 $\Psi\to \Psi+n$，即
$\theta\to\theta+2\pi n$，把 $B$ 平移 $-n\omega_I$，其中
$\omega_I/2\pi$ 是 ${\frak L}$ 的第一 Chern 类。）这一观察使
我们能够在复结构 $K$ 中定义一个方便的 $A$-膜。在 $\Psi=0$ 时，
重要的膜之一是支撑在 $\phi=0$ 处、带平凡 Chan-Paton 丛的膜
${\cal B}_0$。在 $\Psi=n$ 时，我们同样可以取支撑在 $\phi=0$ 处的
膜 ${\cal B}_n$，但此时 Chan-Paton 丛为 ${{{\frak L}}}^n$，并配以
其自然联络。它满足 \eqref{eq:subol}，因而是复结构 $K$ 中的
$A$-膜。

既然设 $\Psi=n$，我们就写 ${\cal B}_{c.c.}^n$ 而不用
${\cal B}_{c.c.}^\Psi$。现在我们来量子化 $({\cal
B}_{c.c.}^n,{\cal B}_n)$ 弦。一切都和从前一样，只是必须把膜
${\cal B}_n$ 的 Chan-Paton 丛包括进来。因此物理态现在是
$K_{\EUM}^{1/2}\otimes {\frak
L}^n\otimes\left(\oplus_{j=0}^n\Omega^{0,j}(\M)\right)$ 的截面。它是
${\cal D}_{K_{\EUM}^{1/2}\otimes {{{\frak L}}}^n}$ 的一个模，而不是
对任何其他 $m$ 的
${\cal D}_{K_{\EUM}^{1/2}\otimes {{{\frak L}}}^m}$ 的模。所以
$({\cal B}_{c.c.}^n,{\cal B}_{c.c.}^n)$ 弦的代数层恰好是
${\cal D}_{K_{\EUM}^{1/2}\otimes {{{\frak L}}}^n}$。

如果 $\Psi$ 不是整数，会发生什么？我们不能指望对一般的 $\Psi$
找到这样一个膜，它给出的模由
$K_{\EUM}^{1/2}\otimes {{{\frak L}}}^\Psi$ 的截面组成，因为除非
$\Psi$ 是整数，这样的模并不存在。然而，我们知道任何 $\Psi$ 下
c.c. 弦的代数都是作用在 ${K_{\EUM}^{1/2}}\otimes
{\frak L}^{f(\Psi)}$ 上的微分算子层，其中 $f(\Psi)$ 是某个函数。
于是，如果我们能证明 $f(\Psi)$ 是线性函数，那么这一点连同
「整数 $\Psi$ 时 $f(\Psi)=\Psi$」这一事实，就蕴含对一切 $\Psi$
都有 $f(\Psi)=\Psi$。

为说明 $f(\Psi)$ 是线性的，考虑 sigma 模型的欧氏作用量：% 式 (11.38)
\setcounter{equation}{37}
\begin{equation}\label{eq:cumber}
S=\int_\Sigma\Phi^*
(\omega-iF-iB)=-i\,\Im\tau \int_\Sigma \Phi^*\left(\cos{\eurm q}\,
\omega_J+\sin{\eurm q}\, \omega_I + i\omega_K\right).
\end{equation}
 （我们略去了形如 $\{Q,V\}$ 的项。）
 当 $\Psi=0$ 时，${\eurm q}=0$，被积式中的二形式在
$T^*\cM(G,C)$ 上成为恰当二形式。相应的一形式势 $\varpi$ 差一个
数值因子就是典范一形式 $p_\alpha\, dq^\alpha$。当
$\Psi\neq 0$ 时，我们能做的最好是把此二形式写成
$T^*\cM(G,C)$ 上的恰当形式加上从 $\cM(G,C)$ 拉回的闭二形式
之和：% 式 (11.39)
\setcounter{equation}{38}
\begin{equation}\label{eq:hexo}
\Omega'={{\Im\tau}\over
\pi}\delta\int_C \,\Tr\left(\phi_z \delta A'_\bz\right)-
i\,\Re\tau \int_C\, \Tr\left(\delta A'_z\delta A'_\bz\right),
\end{equation}

这里我们把 $A'_z$ 定义为 $A'_\bz$ 的 Hermite 共轭的负值。此式中
的第二项是 $\Re\tau$ 乘以 $\M(G,C)$ 上的一个 $(1,1)$-形式，它闭
但不恰当。它是自然线丛 ${\frak L}\to\M$ 的曲率的一个倍数。为
计算转移函数 \eqref{eq:hecox}，我们用传播子展开微扰论，该传播子
是形式 $\Omega'$ 的逆。\eqref{eq:hecox}\ 是由对易子
$[\hat p_\alpha,\hat p_\beta]$ 推出的，它正比于 $\hbar^2$ 或
${1/ ({\rm Im}\,\tau)^2}$。因此我们需要把微扰论做到
$1/({\rm Im}\,\tau)^2$ 阶。在这一阶上，微扰论线性地依赖于
$\Psi={\rm Re}\,\tau$，因为 $\Omega'$ 的逆可以按
${\rm Re}\,\tau$ 与 ${\rm Im}\,\tau$ 展开为级数，其形式概略为
$\sum_{k\geq 0}a_k({\rm Re}\,\tau)^k/({\rm Im}\,\tau)^{k+1}$。
论证到此完成。

$\cM(G,C)$ 上的典范线丛同构于 ${\frak L}^{-2h}$，其中 $h$ 是 $G$
的对偶 Coxeter 数 \cite{ref141}。因此我们也可以说，在任意给定的
$\Psi$ 值下，c.c. 弦的代数都是线丛 ${\frak L}^{\Psi-h}$ 上全纯
微分算子的层。特别地，若令 $\Psi=h$，c.c. 弦的代数就是普通
微分算子层，即作用在函数上的微分算子。由此可见，把 $S$-对偶与
平移 $T^h$（它把 $\Psi$ 平移 $h$）组合起来，我们就能把
$\MH(\LG,C)$ 上的零膜映为 $\cM(G,C)$ 上普通的（未扭化的）
${\cal D}$-模。这是几何 Langlands 对偶更为标准的表述。对某些
$G$，它也是更精确的陈述，只要仔细考虑 \cite{ref5,ref26} 中关于
$C$ 的自旋结构的依赖关系，就能看出这一点。

\begin{yizhu}[易错点]
按：式 (11.26) 把典范参数 $\Psi$ 与 $\theta$ 角、复耦合 $\tau$
联系起来，这里有三套记号容易混淆：$\tau=\theta/2\pi+4\pi
i/e^2$，故 $\Re\tau=\theta/2\pi=\Psi$；而 $t$ 是扭化参数，本节
开头取 $t=1$，之后才允许 $\Psi$ 变动。原文紧接着说「$S$-对偶
一般会变换 ${\rm Re}\,\tau=\theta$」，其中把 $\theta$ 直接写成
${\rm Re}\,\tau$ 是原文的略写，严格地说是 $\theta=2\pi{\rm
Re}\,\tau$；照录于此，读者按上式换算即可。另需注意本节固定
$t=1$、让 $\theta$ 变动，这与第~\ref{sec:canonpar}~节用 $\Psi$
统一参数化的做法一致。
\end{yizhu}

\begin{yizhu}[物理图像]
按：$B$-场 (11.27) 正比于 $\omega_I$，而 Hitchin 纤维在复结构
$I$ 下是复环面、$\omega_I$ 在其上非退化，所以 $B$ 在纤维上的
限制不为零。这正是「一般 $\Psi$ 下 Hitchin 纤维不再是有效
$A$-膜」的根源：包裹该纤维的膜须满足 $F+B|_Y=0$，而
$[B|_Y/2\pi]$ 非整时这对任何 $U(1)$ 联络都无解。物理上可以
说，$\Psi\neq 0$ 时穿透 Hitchin 纤维的磁通量把原来在 $\Psi=0$
处可用的纤维膜「挤掉」了；真正的膜数据须由对偶纤维上的平坦
$U(1)$ 丛给出（参见 Witten 2008 综述关于 Hitchin 纤维化与
$T$-对偶的讨论）。
\end{yizhu}

\begin{yizhu}[易错点]
按：$\lambda$-联络 (11.37) 的定义中，$\lambda$ 只作用在 $\partial
f$ 上（不是 $\bar\partial f$），因此 $\lambda=1$ 给出 $E$ 上的
$\partial$ 算子、$\lambda=0$ 给出 $\End(E)\otimes T^*X$ 的截面。
本节取 $\lambda=i\tan({\eurm q}/2)$；当 $\Psi=0$ 即 ${\eurm q}=0$
时 $\lambda=0$，$D'_z$ 就退化为 Higgs 场 $\phi_z$ 本身。另注意
(11.32) 脚注里的 $w=i\tan({\eurm q}/2)$ 与这里的 $\lambda$ 数值
相同，但二者含义不同：$w$ 参数化复结构 $I_w$，$\lambda$
参数化联络，切勿混为一谈。
\end{yizhu}

\begin{yizhu}[含糊点拨]
按：$f(\Psi)$ 的线性性并非显然，论证结构值得记住。关键量是对易子
$[\hat p_\alpha,\hat p_\beta]\propto\hbar^2\propto 1/({\rm
Im}\,\tau)^2$，故转移函数 (11.39) 只需微扰到 $1/({\rm
Im}\,\tau)^2$ 阶；而在 $\Omega'^{-1}$ 的展开中，${\rm Re}\,\tau$
每出现一次就伴随一个 $1/{\rm Im}\,\tau$ 的因子，于是该阶上的结果
线性依赖于 $\Psi={\rm Re}\,\tau$。再结合整数点 $\Psi=n$ 处已
精确算出的 $f(n)=n$，便得 $f(\Psi)=\Psi$ 对一切 $\Psi$ 成立。
这里「先用微扰论定出函数形式，再用可精确计算的点固定常数」是
本节最有代表性的论证方式。
\end{yizhu}

\begin{yizhu}[物理图像]
按：末段给出本节的总结性图像：一般 $\Psi$ 下 c.c. 弦的代数就是
线丛 ${\frak L}^{\Psi-h}$ 上全纯微分算子的层；当 $\Psi=h$（$h$
为对偶 Coxeter 数）时指数为零，退化为作用在函数上的普通微分
算子。再复合 $T^h$ 平移，就把 $\MH(\LG,C)$ 上的零膜化成
$\cM(G,C)$ 上未扭化的 ${\cal D}$-模，这正是几何 Langlands 对偶
通常的陈述形式。换言之，$\Psi$ 扮演「扭化量」的角色，而 $h$ 是
使扭化恰好被典范线丛的权抵消的临界值；这与正文在 $\Psi=0$ 用
${\cal B}_0$、在 $\Psi=n$ 用 ${\cal B}_n$ 的构造相互印证。
\end{yizhu}

\begin{yizhu}[原文疑误]
按：本片有两处原文笔误，均照录：(1) 讨论 $({\cal
B}_{c.c}^\Psi,{\cal B}_{c.c.}^\Psi)$ 弦的代数结构的那一句中，
左端写作 ${\cal B}_{c.c}^\Psi$（少一个点），按上下文应为
${\cal B}_{c.c.}^\Psi$；(2) 同段把开弦层的下标先写作
$K_{\EUM}^{1/2}$、随后又写作 $K_\M^{1/2}$，$\EUM$ 与 $\M$ 是本文
中同一模空间 $\cM(G,C)$ 的不同宏写法，非笔误但极易误读。另外
(11.31) 下方的文字「${\Im\,\tau}$ 可以任意大，${\eurm q}$ 极其
小」中，${\eurm q}$ 由 $\sin{\eurm q}=-{\Re\,\tau}/{\Im\,\tau}$
确定，需 $\Re\tau$ 有界、$\Im\tau$ 很大时才确实很小。
\end{yizhu}
